\documentclass[conference]{IEEEtran}
\usepackage[utf8]{inputenc}
\usepackage[T1]{fontenc}
\usepackage{cite}
\usepackage{amsmath,amssymb}
\usepackage{graphicx}
\usepackage{booktabs}
\usepackage{url}

\title{Can Deterministic Network Verification Reduce Undetected Faults in LLM‑Generated Configurations?}

\author{
\IEEEauthorblockN{Autonomous AI Researcher}
\IEEEauthorblockA{CNSM 2026 Agentic AI Researcher Track}
}

\begin{document}

\maketitle

\begin{abstract}
We preregistered and executed a paired confirmatory study (study id c1-gnr-vv-2026-0001) to measure whether inserting a deterministic network verifier between an LLM intent\ensuremath{\rightarrow}configuration generator and an emulated execution reduces the proportion of configurations that would execute with operational faults. Using shared\_initial\_candidate semantics (one identical LLM candidate per intent evaluated both without and with verification), we evaluated 72 preregistered paired intents with gpt-5-mini and a canonical deterministic verifier (deterministic\_netops\_validator\_v5). Execution completed with 144 episodes (72 pairs) and 72 model calls. All baseline executions produced no emulator-observed faults (n\_11 = 72, n\_10 = n\_01 = n\_00 = 0). The paired difference in undetected-fault proportion is 0.0 (paired bootstrap 95\% CI [0.0, 0.0], exact McNemar two-sided p = 1.0; bootstrap seed = 1729, resamples = 10000). Under the preregistered shared\_initial\_candidate contract this degenerate confirmatory outcome is expected when identical generated candidates produce no baseline execution faults. We report the preregistration rationale, deterministic execution accounting, representative validator diagnostics, exact inferential outputs, and artifact-grounded recommendations for follow-on confirmatory designs able to detect verifier utility under realistic baseline-error regimes.
\end{abstract}

\section{Introduction}
Generative large language models (LLMs) are increasingly proposed to translate high-level operator intents into device-level network configurations. This intent\ensuremath{\rightarrow}configuration automation can speed change pipelines but risks silently violating reachability, access-control, ordering, or workflow invariants and thereby causing outages if generated text is executed without additional checks. A standard operational mitigation is a generate--verify--repair (GVR) architecture that combines stochastic LLM generation with deterministic verification and, if needed, repair or human review. However, despite many architectural proposals and prototypes, systematic preregistered experiments that measure a verifier's detection performance against executed outcomes under tightly controlled paired designs are scarce.

This paper answers a narrowly scoped, preregistered research question: does inserting a deterministic network verifier between an LLM generator and emulator execution materially reduce the proportion of configurations that execute with operational faults when the identical generated candidate is evaluated both without and with verification? That estimand isolates verifier detection from improvements that could arise from re-sampling, prompt-engineering, or repair-enabled iterations. We executed study c1-gnr-vv-2026-0001 using gpt-5-mini and deterministic\_netops\_validator\_v5 across 72 preregistered paired intents under shared\_initial\_candidate semantics. The run completed with no technical failures, produced a degenerate confirmatory outcome (no baseline emulator faults), and yields a paired difference of 0.0 (95\% CI [0.0,0.0], exact McNemar p = 1.0). Rather than treating this as a null failure, we analyze why it occurred given the preregistered contract and archived artifacts, and we provide concrete, artifact-grounded guidance for designs that can detect verifier utility when baseline error prevalence is non-negligible.

\section{Related Work}
Prior work shows LLMs can perform intent\ensuremath{\rightarrow}configuration translation in constrained vendor-dialect settings and emulator benchmarks [1]. Multi-agent and tool-augmented co-pilot frameworks for NetOps propose integrating verification to reduce regressions [2], and generate--verify--repair patterns have empirical precedent in software and code-generation domains where deterministic checkers validate stochastic outputs and drive repair iterations [3]. Constraint-enforcing decoding and integer-programming constrained decoding reduce structured-output violations in other domains [4], suggesting complementary strategies for NetOps generation. Classic deterministic network-verification techniques (header-space, language-based policy checkers) provide precise invariants for reachability and ACL semantics and are natural verifier candidates, but the literature lacks preregistered paired experiments that report verifier true/false detection rates against emulator-executed failures for identical generated candidates [5]. Our study situates itself at this measurement gap by strictly isolating detection under shared\_initial\_candidate semantics and reporting exact execution-accounting artifacts.

\section{Methodology}
Preregistration and estimand. We preregistered study c1-gnr-vv-2026-0001 and froze the design and analysis prior to observing outcomes. The primary estimand is the paired reduction in undetected operational-fault proportion (paired\_success\_rate\_difference\_guarded\_minus\_baseline). The confirmatory hypothesis H1 targeted an absolute paired reduction of 0.20 with two-sided alpha = 0.05 and power >= 0.80; a sample-size heuristic returned \ensuremath{\approx}63 pairs and we preregistered N = 72 pairs (24 per topology-difficulty stratum) to allow margin for deterministic exclusions.

Shared\_initial\_candidate semantics and experimental contract. The execution\_contract enforced shared\_initial\_candidate semantics: for each intent exactly one initial LLM generation was produced and that identical candidate was evaluated in two conditions. Baseline: execute the shared candidate in an isolated emulator and record emulator fault/no-fault. Guarded: run the canonical deterministic verifier (deterministic\_netops\_validator\_v5) on the same candidate; if the verifier flags violations the guarded outcome is recorded as prevented (no execution); if it does not flag, execute the same candidate in the emulator and record the emulation outcome. This paired structure ensures any observed paired difference arises solely from verifier detection preventing an otherwise-executed fault.

Model, adapter, and verifier configuration. The declared model in execution/model\_configuration.json is gpt-5-mini (provider = openai\_responses, max\_output\_tokens = 2000, maximum\_attempts\_per\_call = 1). execution/model\_configuration.json records temperature = null; this indicates provider-default runtime sampling semantics and is archived as a residual sampling-parameter uncertainty. The adapter family was hosted\_netops\_gvr\_v1. The canonical verifier is deterministic\_netops\_validator\_v5 (validator\_version = "5.0") using a full\_intent\_constraint\_validation rule set that outputs per-episode validity verdicts and structured validator\_traces (parsed\_command\_count, required\_command\_count, validation\_before/after states, violation\_codes, difficulty metadata). All verifier decisions and traces were archived per episode under execution/scoring/.

Task-bank construction, contamination controls, and stratification. Tasks were deterministically generated by deterministic\_netops\_task\_generator\_v5. Each task manifest entry encodes a natural-language intent, a machine-structured required\_sequence (ordered field changes), allowed\_touched\_fields, initial and expected final states, and difficulty labels (medium/hard/very\_hard). The preregistration required deterministic exact-match contamination checks against public corpora with explicit exclusion rules. analysis/contamination\_summary.csv records zero exact-match exclusions for confirmatory pairs. The confirmatory sample retained the preregistered stratification (24 pairs per stratum).

Execution protocol, retries, and deterministic accounting. The missingness\_plan allowed up to two additional retries (three attempts total) for transient hosted-API errors on generation calls and required full logging of retries and provider events. The run started at 2026-08-31T06:56:23.065707Z and completed at 2026-08-31T07:06:28.665518Z (\ensuremath{\approx}10.1 minutes wall-clock). The execution manifest reports planned\_episode\_count = 144, completed\_episode\_count = 144, failed\_episode\_count = 0, and model\_calls\_used = 72 (one shared generation per pair). Provider-call audit (deterministic\_reconciliation.json) shows hosted\_model\_call\_event\_count = 72, completed\_hosted\_model\_call\_event\_count = 72, cache\_miss\_count = 72, and stage\_counts = \{shared\_initial\_generation: 72\}.

Scoring, validation, and inferential procedures. Each episode's verifier emits a score and validator\_trace and the primary analysis used paired\_binary\_analysis\_v1 to form the 2\ensuremath{\times}2 paired contingency table (guarded \ensuremath{\times} baseline). The preregistered inferential test is an exact McNemar two-sided test when discordant pairs exist. We computed paired bootstrap-percentile 95\% confidence intervals for the paired difference using 10,000 resamples with bootstrap\_seed = 1729 (analysis/analysis\_log.jsonl). Pre-specified subgroup confirmatory comparisons (topology complexity and policy type) were registered with Holm correction; these are archived but become non-informative in the absence of discordant pairs. All analysis was executed deterministically by the registered analysis executor and archived under analysis/.

Additional methodological notes (artifact-grounded). The validator's full\_intent\_constraint\_validation rule set enforces sequence and field constraints encoded in the task manifests: validator traces include parsed\_command\_count and required\_command\_count and a difficulty.level tag (examples below). Per-episode JSON scoring artifacts under execution/scoring/ record normalized\_configuration, validation\_before/after final\_state snapshots, and violation\_codes arrays that were systematically inspected during analysis.

\section{Results}
Execution and manifest summary. The autonomous pipeline completed without technical failures and met the preregistered accounting: planned\_episode\_count = 144; completed\_episode\_count = 144; failed\_episode\_count = 0; model\_calls\_used = 72. The run-level audit (deterministic\_reconciliation.json) records hosted\_model\_call\_event\_count = 72, completed\_hosted\_model\_call\_event\_count = 72, cache\_miss\_count = 72, and stage\_counts = \{shared\_initial\_generation: 72\}. analysis/condition\_summary.csv reports baseline successes = 72, baseline failures = 0 (success\_rate = 1.0) and guarded successes = 72, guarded failures = 0 (success\_rate = 1.0). analysis/paired\_contingency\_table.csv contains the 2\ensuremath{\times}2 counts used for inference.

Paired contingency, inference, and bootstrap diagnostics. Aggregating across the 72 confirmatory pairs produced contingency counts n\_11 = 72, n\_10 = 0, n\_01 = 0, n\_00 = 0 (guarded \ensuremath{\times} baseline). The paired estimate paired\_success\_rate\_difference\_guarded\_minus\_baseline = 0.0. The preregistered exact McNemar two-sided test returns p = 1.0 given zero discordant pairs. The paired bootstrap-percentile 95\% CI (bootstrap\_seed = 1729; bootstrap\_resamples = 10000) is [0.0, 0.0], produced by resampling the observed paired outcomes and recomputing the paired difference; analysis/analysis\_log.jsonl records these bootstrap parameters. These inferential outputs are archived in analysis/paired\_contingency\_table.csv and analysis/condition\_summary.csv.

Per-episode validator traces and representative exemplars. Every episode produced a deterministic validator\_trace (execution/scoring/*.json). Representative extracts from archived per-episode traces show the validator executed full\_intent\_constraint\_validation checks and produced no violation flags across examples: 
- pair-000001 (task-000001): difficulty.level = "medium"; parsed\_command\_count = 4; required\_command\_count = 4; observed\_touched\_field\_count = 3; validator\_version = "5.0"; violation\_count = 0.
- pair-000002 (task-000002): difficulty.level = "hard"; parsed\_command\_count = 2; required\_command\_count = 2; observed\_touched\_field\_count = 2; violation\_count = 0.
- pair-000003 (task-000003): difficulty.level = "hard"; parsed\_command\_count = 6; required\_command\_count = 6; observed\_touched\_field\_count = 4; violation\_count = 0.
The archived manifest lists representative\_tasks with intents, required\_sequence encodings, and responses (execution/responses/task-00000X-*.txt). Across the full confirmatory set the verifier flagged zero violations (analysis/paired\_contingency\_table.csv and analysis/condition\_summary.csv reflect this aggregate result). The per-episode JSONs under execution/scoring/ provide full validator\_before/after state snapshots for reproducibility.

Stratification and subgroup diagnostics. The preregistered stratification (24 pairs per difficulty stratum: medium/hard/very\_hard) was preserved. Because no discordant pairs occurred, the pre-specified subgroup confirmatory comparisons (topology complexity and policy type with Holm correction) produced no informative contrasts; their computation and null outputs are archived but substantively uninformative under a zero-prevalence baseline. The absence of discordance also rendered sensitivity checks (e.g., treating excluded pairs as failures) numerically identical to the main analysis because analysis/contamination\_summary.csv reports zero exact-match exclusions and analysis/missingness\_summary.csv reports complete\_pairs = 72 and zero failed or unscorable episodes.

Sensitivity and exploratory diagnostics executed. We executed the preregistered sensitivity analyses archived in analysis/: (a) complete\_pair\_primary\_with\_failure\_as\_zero\_sensitivity (treating any excluded pairs as failures) --- not applicable because no pairs were excluded; (b) bootstrap CI stability (10k resamples, seed = 1729) --- reproduced CI [0.0,0.0]; (c) exploratory taxonomy of missed faults --- vacuous here because no missed faults were observed. All analysis logs and artifact files for these diagnostics are archived under analysis/ (see analysis/analysis\_log.jsonl, analysis/missingness\_summary.csv, analysis/contamination\_summary.csv).

Artifact-grounded account of the degenerate outcome. The degenerate confirmatory result (no baseline emulator faults) is traceable in archived artifacts to measurable design choices: (1) the task generator produced tightly constrained intents with explicit required\_sequence fields and allowed\_touched\_fields encoded in each task manifest (representative entries in execution/task\_manifest.jsonl), making the desired configuration sequence unambiguous; (2) the model used for initial generation (declared\_model\_version = "gpt-5-mini" in execution/model\_configuration.json) produced candidates that matched the manifest-encoded required sequences across the synthetic corpus; and (3) the verifier's full\_intent\_constraint\_validation rule set directly enforces those same manifest constraints, so verifier expectations and generation targets were highly congruent. Together these factors explain why the shared generated candidates both passed verification and executed without emulator-observed faults.

Reproducibility artifacts and deterministic checks. Key archived artifacts enabling deterministic reconciliation include: execution/raw\_results.jsonl (raw per-episode records), execution/task\_manifest.jsonl (task payloads and required\_sequence encodings), execution/model\_configuration.json (declared\_model\_version = gpt-5-mini; temperature = null; max\_output\_tokens = 2000), analysis/paired\_contingency\_table.csv, analysis/condition\_summary.csv, deterministic\_reconciliation.json, and per-episode validator traces under execution/scoring/. The bootstrap procedure used seed = 1729 and 10,000 resamples; these values are recorded in analysis/analysis\_log.jsonl. Provider-call accounting and model-call counts are recorded in deterministic\_reconciliation.json to allow bit-for-bit reconciliation of run-level totals.

Summary numerical takeaways. Under the preregistered shared\_initial\_candidate semantics and the chosen synthetic task bank the observed baseline undetected-fault proportion = 0/72 = 0.0, guarded undetected-fault proportion = 0/72 = 0.0, paired difference = 0.0, exact McNemar p = 1.0, and paired bootstrap 95\% CI = [0.0,0.0]. All confirmatory diagnostic artifacts and per-episode validator\_traces that substantiate these numbers are archived and referenced above.

\section{Discussion}
Confirmatory interpretation and boundary conditions. The preregistered paired design with shared\_initial\_candidate semantics validly measures verifier detection on identical candidates; because the identical generated candidates produced zero emulator faults the verified paired outcome is a ceiling/null result: the verifier could not reduce executed faults that did not occur. The numeric outputs (n\_11 = 72; paired difference = 0.0; McNemar p = 1.0; bootstrap CI [0.0,0.0]) are direct consequences of the preregistered contract and observed data and therefore correctly reported as a degenerate confirmatory result.

Artifact-grounded causes of the ceiling outcome. Archived artifacts allow us to identify measurable factors that explain the ceiling outcome: (1) the deterministic task generator produced constrained, high-clarity intents with explicit required\_sequence fields and allowed\_touched\_fields encoded in each task manifest, increasing the likelihood of unambiguous correct candidates; (2) gpt-5-mini (declared\_model\_version in execution/model\_configuration.json) generated sequences that matched those structured constraints across the synthetic corpus; and (3) the verifier rule set (full\_intent\_constraint\_validation in deterministic\_netops\_validator\_v5) aligns closely with the required\_sequence constraints present in the task manifests, making verification and generation objectives strongly congruent. These archived, measurable factors together account for the absence of baseline emulator faults.

Design implications: when verification benefit is measurable. To empirically measure practical verifier benefits, confirmatory contracts must present non-negligible baseline error prevalence or allow verifier-mediated candidate changes. Artifact-grounded, preregisterable alternatives include: (A) independent-generation arms---allow separate initial generations per condition so verifier-triggered repairs or alternative sampling may change executed candidate distributions; (B) repair-enabled flows---permit a deterministic repair call when the verifier flags violations and execute repaired candidates to measure end-to-end improvement; (C) adversarial/fault-injection task banks---seed tasks with ambiguous intents, ordering subtleties, or vendor-edge cases to raise baseline error prevalence; (D) multi-model and sampling sweeps---include multiple model versions and explicit sampling parameter settings (non-null temperatures) to quantify model- and sampling-sensitivity. All these options are compatible with the archived execution and analysis infrastructure and can be preregistered using the same evidence-recording conventions.

Threats to validity revisited (artifact-supported). Internal validity: by design the shared\_initial\_candidate semantics isolates detection but prevents verifier-mediated candidate resampling from influencing outcomes; this tradeoff explains why a degenerate result may occur when baseline error prevalence is low. Construct validity: emulator-observed faults are a conservative proxy for operational harm but do not capture traffic-dependent timing, vendor-specific hardware idiosyncrasies, or live-traffic emergent failures. External validity: experiments used a single declared model (gpt-5-mini), a synthetic task bank, and one deterministic verifier; extrapolation to other models, vendor dialects, or production hardware is limited. Conclusion validity: because baseline error prevalence was zero in this corpus the confirmatory design had no power to detect verifier benefit; the preregistered power target (detect absolute paired reduction 0.20) becomes moot under a zero-prevalence baseline.

Operational implications and measured tradeoffs. The confirmed ceiling outcome does not imply verifiers are unnecessary; rather, it shows that under certain tightly constrained synthesis and evaluation conditions an LLM alone produced valid device sequences on the synthetic corpus. In production settings with ambiguous intents, incomplete constraints, or vendor heterogeneity, verifiers remain a mechanistic guard. Measuring verifier costs (false positives that block safe changes) and benefits (true positives preventing faults) requires task banks with non-negligible baseline failure prevalence; the archived artifacts and preregistration provide a reproducible template for such follow-on evaluations.

Reproducibility notes (artifact pointers and deterministic checks). All execution and analysis artifacts are archived in the evidence bundle. Key artifacts referenced in this paper include: execution/raw\_results.jsonl, execution/task\_manifest.jsonl, execution/model\_configuration.json, analysis/paired\_contingency\_table.csv, analysis/condition\_summary.csv, analysis/analysis\_log.jsonl, deterministic\_reconciliation.json, and per-episode validator traces under execution/scoring/. The analysis bootstrap used seed = 1729 and 10,000 resamples; analysis/analysis\_log.jsonl records these values. The run-level timing, provider-call audit, and artifact counts are recorded in execution/ and analysis/ manifests to enable deterministic reconciliation of the paired accounting. The model configuration records temperature = null (execution/model\_configuration.json), which we report as provider-default sampling semantics (a residual sampling-parameter uncertainty archived in the evidence).

\section{Limitations}
\begin{itemize}
\item Confirmatory ceiling/null outcome: the baseline generator produced zero emulator faults on the preregistered task set, preventing direct measurement of verifier detection benefit.
\item Shared\_initial\_candidate semantics: the design measures verifier effect on the same generated candidate only and does not assess independent per-condition generation, multi-sample generation, or repairs unless explicitly permitted by the contract.
\item Single-model scope: experiments used one declared model (gpt-5-mini); results do not imply generality across models or versions.
\item Synthetic/emulated tasks: task corpus and emulator executions differ from production devices and live traffic; external validity to production deployments is limited.
\item Sampling-parameter uncertainty: execution/model\_configuration.json shows temperature = null (sampling parameters unset); null indicates provider-default runtime sampling semantics and is a residual uncertainty recorded in the archived configuration.
\item Residual contamination risk: deterministic provenance and exact-match checks were applied, but private training-data contamination of the hosted model cannot be fully ruled out and remains residual uncertainty.
\end{itemize}

\section{Conclusion}
We executed a preregistered paired evaluation (c1-gnr-vv-2026-0001) under shared\_initial\_candidate semantics to test whether a canonical deterministic verifier reduces the proportion of LLM-generated configurations that would execute with operational faults. For the chosen model (gpt-5-mini), verifier (deterministic\_netops\_validator\_v5), and synthetic/emulated task bank, baseline executions produced zero emulator faults and the verifier did not change outcomes (paired difference = 0.0; bootstrap 95\% CI [0.0,0.0]; exact McNemar p = 1.0). This artifact-grounded null/ceiling outcome is internally consistent with the preregistered design: when identical generated candidates produce no executed faults a verifier cannot reduce executed faults under the shared\_initial\_candidate contract. Measuring verifier utility under realistic error regimes requires preregistered contracts that permit independent or multiple generator samples, repair-enabled flows, adversarial task sets, and multi-model comparisons. The archived artifacts (responses, task manifests, validator traces, execution manifest, and analysis logs) provide a reproducible baseline and a template for those follow-on confirmatory designs and power calculations.

\section*{Disclosure Statement}
Autonomy and artifacts: This study was executed autonomously after an autonomy lock. Human role: a Research Director selected the topic family and approved the immutable initial master prompt prior to the autonomy lock; no human manually edited technical manuscript content after the lock. Models and tools: initial generation used gpt-5-mini (declared\_model\_version recorded in execution/model\_configuration.json) orchestrated via the hosted\_netops\_gvr\_v1 adapter; deterministic\_netops\_validator\_v5 (validator\_version = "5.0") served as the canonical verifier; an isolated emulator executed candidate configurations. Preregistration: study c1-gnr-vv-2026-0001 was preregistered and frozen before observing outcomes. Immutable master prompt reference: provenance/master\_prompt.txt (SHA-256: 1872df1e\allowbreak{}1805d2d9\allowbreak{}6940456c\allowbreak{}a016bd66\allowbreak{}5d1d5196\allowbreak{}add77f5a\allowbreak{}cdf1582b\allowbreak{}b39b15ba). Key archived artifacts include execution/raw\_results.jsonl, execution/task\_manifest.jsonl, execution/model\_configuration.json, analysis/paired\_contingency\_table.csv, analysis/condition\_summary.csv, deterministic\_reconciliation.json, and per-episode validator traces under execution/scoring/. The model configuration records temperature = null (provider-default sampling semantics archived in execution/model\_configuration.json). Provider-call audit: hosted\_model\_call\_event\_count = 72. No human manual manuscript editing or content insertion occurred after the autonomy lock.

\begin{thebibliography}{99}
\bibitem{ref1} Nguyen Tu, Sukhyun Nam, James Won‐Ki Hong, "Intent‐Based Network Configuration Using Large Language Models", 2024, 10.1002/nem.2313
\bibitem{ref2} Zhaodong Wang, Samuel Lin, Guanqing Yan, Soudeh Ghorbani, Minlan Yu, Jiawei Zhou, Nathan Hu, Lopa Baruah, Sam Peters, Srikanth Kamath, Jian Yang, Ying Zhang, "Intent-Driven Network Management with Multi-Agent LLMs: The Confucius Framework", 2025, 10.1145/3718958.3750537
\bibitem{ref3} Rafael Cadenas, "Convergent Correctness in Stochastic Code Generation: A Generate-Verify-Repair Architecture for Deterministic Validation of Autonomous AI Coding Agents in Regulated Environments", 2026, 10.2139/ssrn.6754899
\bibitem{ref4} http://citeseerx.ist.psu.edu/viewdoc/summary?doi=10.1.1.300.8659
\end{thebibliography}

\end{document}
